% De acuerdo con la NTC 1486, la introducción presenta la importancia, el origen, los objetivos, el alcance, las limitaciones y la metodología del trabajo, sin anticipar resultados ni conclusiones (2 a 4 páginas).

Buscar empleo remoto exige recorrer varias plataformas, y muchas de ellas buscan comparando las palabras de la consulta con las de la oferta. Ese mecanismo funciona cuando quien busca y quien publica describen el cargo con los mismos términos, y falla cuando no lo hacen: \enquote{Backend Developer} y \enquote{Software Engineer}, \enquote{JS} y \enquote{JavaScript}, \enquote{trabajo con datos} y \enquote{Analytics Engineer} pueden referirse a lo mismo sin compartir una sola palabra. La recuperación de información llama a este problema \emph{desajuste de vocabulario}. Las grandes plataformas de reclutamiento lo abordan con representaciones semánticas aprendidas de sus propios datos de interacción \citep{HaThuc2017,Shen2024}, un recurso que no está al alcance de un agregador pequeño.

\textit{Retrieval-Augmented Generation} (RAG), o generación aumentada por recuperación \citep{Lewis2020}, ofrece una alternativa construida con componentes disponibles: recupera los fragmentos más cercanos a la consulta y le pide a un modelo de lenguaje que redacte una respuesta a partir de ellos. Sobre ese esquema se han propuesto mejoras como los filtros sobre los atributos de la oferta, la búsqueda híbrida \citep{Cormack2009}, el reordenamiento de los candidatos con un modelo más costoso \citep{Nogueira2019} y los modelos de \textit{embeddings} multilingües \citep{Wang2024mE5}. Cada una promete mejorar la recuperación y cada una añade latencia y complejidad, pero en la búsqueda de empleo no se sabe cuánto aporta cada una frente a una línea base léxica fuerte. Usar el sistema tampoco basta para saberlo, porque un pipeline RAG responde con fluidez aunque su recuperación o sus datos fallen \citep{Barnett2024}.

De ahí surge la pregunta que orienta este trabajo: \emph{¿en qué medida un pipeline RAG mejora la recuperación de ofertas de empleo frente a la búsqueda por palabras clave, qué componentes del pipeline explican esa mejora y con qué costo en latencia?} Para responderla se toma como caso de estudio JobAI, una plataforma que agrega ofertas de tres fuentes públicas orientadas al empleo remoto y responde consultas en lenguaje natural mediante un pipeline RAG. El autor la desarrolló antes de esta tesis, sin un protocolo de evaluación formal, y por eso ofrece un caso real en el que medir lo que cada componente aporta.

El objetivo general es evaluar en qué medida un pipeline RAG aplicado a la búsqueda de empleo remoto mejora la calidad de la recuperación frente a la búsqueda por palabras clave, e identificar el aporte y el costo de cada uno de sus componentes. Para alcanzarlo, el trabajo parte de la literatura sobre los límites de la búsqueda léxica y los componentes que RAG propone para superarlos, diseña con ella un protocolo de evaluación repetible, lo aplica para mejorar JobAI en dos iteraciones y, con la configuración resultante, cuantifica el aporte de cada componente, la fidelidad de las respuestas y el costo en latencia. El Capítulo~\ref{cap:problema} enuncia los cuatro objetivos específicos que siguen ese recorrido.

Lo que se estudia es el pipeline de consulta de JobAI, que se mejora en dos iteraciones y se evalúa con un protocolo diseñado para este trabajo: un corpus real de ofertas congelado para poder repetir los experimentos, consultas de prueba en inglés y en español, BM25 como línea base con y sin los mismos filtros que el pipeline, una escalera que agrega los componentes uno a uno y juicios de relevancia de un modelo de lenguaje validado contra juicios humanos, con un conjunto de desarrollo separado para las decisiones de configuración. La primera iteración corrige el sistema para que haga lo que su documentación describe, y la segunda incorpora los componentes cuyo aporte se mide. El código, el corpus congelado, las consultas, los juicios y los resultados de cada corrida están en el repositorio \url{https://github.com/esteban0406/RAG-JobPosting} (Anexo~\ref{anx:repositorio}).

El estudio evalúa la calidad técnica del pipeline sobre ofertas en inglés, mayoritariamente del sector tecnológico, y trata cada consulta de forma aislada, sin usuarios reales; la Sección~\ref{sec:alcance} detalla lo que queda fuera. Dentro de ese alcance, el trabajo aporta una evaluación por componentes de un pipeline RAG para la búsqueda de empleo frente a BM25, que separa el efecto de cada componente y su costo en latencia; un registro de las fallas que aparecieron al evaluar un sistema RAG en funcionamiento, con el origen de cada detección y lo que las métricas de RAGAS sobre la respuesta final alcanzaron a ver de ellas, y, como resultado secundario, el costo de la búsqueda según el motor de almacenamiento vectorial, a la escala del corpus real y a escalas mayores.

El documento sigue el orden en que se construyó la respuesta. El Capítulo~\ref{cap:problema} plantea el problema, los objetivos y el marco conceptual, y el Capítulo~\ref{cap:estado-arte} revisa la literatura sobre los componentes del pipeline y su evaluación. El Capítulo~\ref{cap:metodologia} describe el protocolo y el Capítulo~\ref{cap:artefacto}, la arquitectura de JobAI y sus dos iteraciones. El Capítulo~\ref{cap:resultados} presenta y discute los resultados, incluidas las fallas encontradas, y el Capítulo~\ref{cap:conclusiones} responde a cada objetivo y propone trabajo futuro. Los anexos reúnen los detalles del protocolo y de JobAI y un mapa del repositorio.
